\begin{abstract}
As computing architectures trend toward heterogeneous integration, Dynamic Thermal Management (DTM) and temperature-aware scheduling have become critical research domains. However, recent literature indicates a growing reliance on external thermal solvers (e.g., HotSpot, 3D-ICE, FEM) due to pervasive skepticism regarding the physical fidelity of native architectural thermal models. In this paper, we systematically investigate the thermal infrastructure of the gem5 architectural simulator and expose a latent absolute-zero (\SI{0}{\kelvin}) initialization anomaly. We demonstrate mathematically how this artifact induces unphysical junction cooling during the initial simulation epochs, significantly distorting transient thermal behavior and potentially misleading DTM evaluations. To restore physical admissibility, we not only propose an architectural patch but also introduce a comprehensive five-way validation framework encompassing closed-form analytical solutions, the buggy native simulation, the patched simulation, an independent Python-based RC solver, and industry-standard SPICE circuit simulations. Our results corroborate that the patched native thermal model achieves a transient and steady-state error margin of less than \SI{0.08}{\degreeCelsius} compared to SPICE baselines. This work bridges a critical methodological gap, providing the much-needed physical admissibility proof for native architectural thermal simulations.
\end{abstract}
